\noindent This section repeats the result of Secs.~6 and 7 (after the exact cancellations and the symmetrisation of the charges), Group by Group and Row by Row, with (a) every colour structure written as a sum of products of \emph{differences} of Wilson lines, as in the LO structure $\big[U^{ce'}(y')-U^{ce'}(x')\big]\big[U^{ce}(y)-U^{ce}(x)\big]$, and (b) every longitudinal integral written with the momentum fraction $\xi$. Everything else (prefactors, transverse kernels, phases, the order of the charges, the c.c.) is that of Secs.~6 and 7. Every rewriting is an exact identity (any $N_{c}$, nothing dropped; checks in Sec.~\ref{sec:F3}). The conventions of Sec.~\ref{sec:conv2} apply, in particular: principal value in Group I Row III, pure glue, and the terms that must be integrated together.

\paragraph*{Longitudinal momentum fraction.} Two definitions (author, 2026-11-23):
\begin{itemize}
\item Group II Row IV (as for DGLAP, Sec.~3): $\xi=\frac{k^{+}}{p^{+}+k^{+}}$, the fraction of the parent momentum carried by the observed gluon ($k^{+}/\xi$ is the parent momentum).
\item All other rows: $\xi=\frac{p^{+}}{k^{+}}$.
\end{itemize}
Jacobians and limits: Group I: $dp^{+}=k^{+}d\xi$, $0\le\xi\le1$; Group III: $dp^{+}=k^{+}d\xi$, $0\le\xi<\infty$; Group II Row IV: $p^{+}=k^{+}\frac{1-\xi}{\xi}$, $dp^{+}=-k^{+}\frac{d\xi}{\xi^{2}}$, $p^{+}=0\leftrightarrow\xi=1$, $p^{+}=\vee-k^{+}\leftrightarrow\xi=\frac{k^{+}}{\vee}$, so that
\begin{align}
\int_{0}^{\vee-k^{+}}\frac{dp^{+}}{2\pi}\,\frac{p^{+}}{(p^{+}+k^{+})^{2}}\,e^{-ik\cdot(y-y')}\,e^{-ik\cdot(y-y')\frac{p^{+}}{k^{+}}}=\int_{k^{+}/\vee}^{1}\frac{d\xi}{2\pi}\,\frac{1-\xi}{\xi}\,e^{-ik\cdot(y-y')\frac{1}{\xi}}\,.\notag
\end{align}
The recoil phases become $e^{-ik\cdot(a)\frac{p^{+}}{k^{+}}}=e^{-ik\cdot(a)\,\xi}$. The plus prescriptions become (same subtraction points; the prescription acts on the whole $\xi$-dependent function $f$ of the term):
\begin{align}
\int_{0}^{k^{+}}dp^{+}\Big[\frac{1}{p^{+}}\Big]_{+}f&=\int_{0}^{1}d\xi\Big[\frac{1}{\xi}\Big]_{+}f\,,\qquad \int_{0}^{k^{+}}dp^{+}\Big[\frac{k^{+}}{p^{+}}\Big]_{+}f=k^{+}\int_{0}^{1}d\xi\Big[\frac{1}{\xi}\Big]_{+}f\,,\qquad \int_{0}^{1}d\xi\Big[\frac{1}{\xi}\Big]_{+}f(\xi)=\int_{0}^{1}\frac{d\xi}{\xi}\big[f(\xi)-f(0)\big]\,,\notag\\
\int_{0}^{k^{+}}dp^{+}\Big[\frac{k^{+}}{k^{+}-p^{+}}\Big]_{+}f&=k^{+}\int_{0}^{1}d\xi\Big[\frac{1}{1-\xi}\Big]_{+}f\,,\qquad \int_{0}^{1}d\xi\Big[\frac{1}{1-\xi}\Big]_{+}f(\xi)=\int_{0}^{1}\frac{d\xi}{1-\xi}\big[f(\xi)-f(1)\big]\,,\notag\\
\int_{0}^{\infty}dp^{+}\Big[\frac{k^{+}}{p^{+}(p^{+}+k^{+})}\Big]_{+}f&=\int_{0}^{\infty}d\xi\Big[\frac{1}{\xi(1+\xi)}\Big]_{+}f\,,\qquad \int_{0}^{\infty}d\xi\Big[\frac{1}{\xi(1+\xi)}\Big]_{+}f(\xi)=\int_{0}^{\infty}\frac{d\xi}{\xi(1+\xi)}\big[f(\xi)-f(0)\big]\,,\notag
\end{align}
and the instantaneous terms of Group III: $\int_{0}^{\infty}d\xi\,e^{-ik\cdot(z-y')\,\xi}=\pi\delta(s)-i\,\mathcal{P}\frac{1}{s}$, $s=k\cdot(z-y')$ (with the conditions of Sec.~\ref{sec:conv2}). The energy denominators become $\frac{1}{p^{+}+k^{+}}=\frac{1}{k^{+}(1+\xi)}$, $\frac{k^{+}}{p^{+}(y-x)^{2}+(k^{+}-p^{+})(y-z)^{2}}=\frac{1}{\xi(y-x)^{2}+(1-\xi)(y-z)^{2}}$, $\frac{1}{p^{+}(x-z)^{2}+k^{+}(x-w)^{2}}=\frac{1}{k^{+}}\frac{1}{\xi(x-z)^{2}+(x-w)^{2}}$; the factors $k^{+}$ are collected in the prefactors (every row has the prefactor $\frac{1}{k^{+}}$).

\paragraph*{Colour structures.} The identities used (exact, any $N_{c}$): $f^{gde}U^{db}(p)U^{ec}(p)=U^{gh}(p)f^{hbc}$ (invariance of $f$), $U^{mh}(p)U^{kh}(p)=\delta^{mk}$, $f^{acd}f^{bcd}=N_{c}\delta^{ab}$; e.g.\ $f^{hbc}U^{gh}(w)-f^{gde}U^{db}(x)U^{ec}(y)=-f^{gde}\big\{[U^{db}(x)-U^{db}(w)]U^{ec}(y)+U^{db}(w)[U^{ec}(y)-U^{ec}(w)]\big\}$. In the singlet form the contractions $C^{abc}d^{abc}$ (three-$\rho$ part) and $C^{abc}f^{abc}$ at coinciding charges (contact part) of a product of three matrices $A,B,C$ (each a product of Wilson lines and differences, one of them, $O$, a Wilson line or a product of Wilson lines) are written as
\begin{align}
f^{ijk}d^{lmn}O^{il}B^{jm}C^{kn}=i\,{\rm Tr}\big\{T^{h}O^{\dagger}C\,D^{h}B^{\dagger}O\big\}\,,\qquad f^{ijk}f^{lmn}O^{il}B^{jm}C^{kn}={\rm Tr}\big\{T^{e}O^{\dagger}C\,T^{e}B^{\dagger}O\big\}\,,\notag
\end{align}
with $B$, $C$ the differences, e.g.\ ${\rm Tr}\big\{T^{h}U^{\dagger}(x)\big(U(z)-U(y)\big)D^{h}\big(U^{\dagger}(x')-U^{\dagger}(w')\big)U(x)\big\}$. Each row's singlet equation has the form of Sec.~7 with $\frac{N_{c}}{(N_{c}^{2}-4)(N_{c}^{2}-1)}(t\,i)\,\sigma\,[\dots]$, $t=2$ for the rows with an anticommutator and $t=1$ for the ordered products, and the contact coefficients $(\pm i)$, $(\pm\frac{i}{2})$ of Sec.~7.

\fixnote{\textbf{What cannot be written as a product of differences only.} (1) Every three-$\rho$ colour tensor $C^{abc}$ has three adjoint indices, one for each charge. In every term exactly two of the three colour legs carry a difference of Wilson lines; the third leg carries one Wilson line $U(p)$ (or a product $U^{\dagger}(z)U(\dots)$, or the unit matrix). This third Wilson line cannot be turned into a difference: writing $U(p)=[U(p)-U(q)]+U(q)$ adds a term that does not vanish. Therefore the three-$\rho$ structures are given as ``$[U-U]\times[U-U]\times U\times f$'' and, in singlet form, as traces with two differences and one Wilson line ($O$ above). (2) In Group II Row I, Group III Rows II and III (second part) and in the contact terms of Group II Row I and Group III Row II the factorisation produces products $U^{\dagger}(z)\,[U(z)-U(y)]$ and $[U^{\dagger}(w)-U^{\dagger}(z)]\,U(x')$, i.e.\ a difference times a Wilson line in one leg; this is the closest exact form (the same reason). (3) In the three-$\rho$ part of Group I Row II the two traces of the factorised form add up to a shorter expression with one difference (the term $f^{hbc}U^{gh}(w)$ of the tensor drops out of $C^{abc}d^{abc}$ because $f^{hbc}d^{abc}=0$, Sec.~7); the factorised form is kept for uniformity, both are exact. (4) The genuine two-$\rho$ structures and all contact (two-$\rho$) structures are sums of products of two differences, with $f$, $N_{c}$ and one Wilson line; nothing is left unfactorised.}
